\documentclass[preview]{standalone}
\usepackage{amsmath}
\usepackage{amssymb}
\usepackage{stellar}
\usepackage{definitions}
\usepackage{bettelini}
\begin{document}
\id{diffeq-variation-constants}
\genpage
\section{Nonhomogeneous equations}

\begin{snippetproposition}{diffeq-nonhomogeneous-affine-space}{General solution of a nonhomogeneous equation}
    Let \(L[y]=y^{(k)}+\sum_{j=0}^{k-1}a_jy^{(j)}\), with real constant
    coefficients, and let \(b\in C(I)\) on an open interval \(I\).
    If \(y_*\) solves \(L[y_*]=b\) and \(w_1,\ldots,w_k\) form a basis
    of the homogeneous solution space on \(I\), then all solutions of \(L[y]=b\) are
    \[
        y=y_*+\sum_{j=1}^k c_jw_j,\qquad c_j\in\realnumbers.
    \]
\end{snippetproposition}

\begin{snippetproof}{diffeq-nonhomogeneous-affine-space-proof}{diffeq-nonhomogeneous-affine-space}{General solution of a nonhomogeneous equation}
    If \(L[w]=0\), linearity gives \(L[y_*+w]=b+0=b\).
    Conversely, if \(L[y]=L[y_*]=b\), then \(L[y-y_*]=b-b=0\).
    Expand this difference in the homogeneous basis.
\end{snippetproof}

\section{Wronskians and variation of constants}

\begin{snippetdefinition}{diffeq-wronskian-definition}{Wronskian matrix}
    For \function[functions] \(w_1,\ldots,w_k\in C^{k-1}(I)\), their
    \emph{Wronskian matrix} and \emph{Wronskian determinant} are
    \[
        W(x)=
        \begin{pmatrix}
        w_1&\cdots&w_k\\
        w_1'&\cdots&w_k'\\
        \vdots&\ddots&\vdots\\
        w_1^{(k-1)}&\cdots&w_k^{(k-1)}
        \end{pmatrix}(x),
        \qquad \mathcal W(x)=\det W(x).
    \]
\end{snippetdefinition}

\begin{snippetproposition}{diffeq-fundamental-wronskian-nonzero}{Nonvanishing of a fundamental Wronskian}
    Let \(w_1,\ldots,w_k\) be linearly independent solutions on an open interval
    \(I\) of \(y^{(k)}+\sum_{j=0}^{k-1}a_jy^{(j)}=0\), with constant real
    coefficients. Their Wronskian matrix is invertible at every \(x\in I\).
\end{snippetproposition}

\begin{snippetproof}{diffeq-fundamental-wronskian-nonzero-proof}{diffeq-fundamental-wronskian-nonzero}{Nonvanishing of a fundamental Wronskian}
    If \(W(x_0)c=0\) for a nonzero vector \(c\), the solution
    \(w=\sum c_jw_j\) has all \(k\) initial data zero at \(x_0\).
    Uniqueness forces \(w=0\), contradicting independence.
    This argument uses that the \function[functions] solve the same equation;
    independence alone does not guarantee a nonzero Wronskian for arbitrary \function[functions].
\end{snippetproof}

\begin{snippettheorem}{diffeq-variation-constants-theorem}{Variation of constants}
    Let \(L[y]=y^{(k)}+\sum_{j=0}^{k-1}a_jy^{(j)}\), with real constant
    coefficients, and let \(b\in C(I)\) on an open interval \(I\).
    Let \(w_1,\ldots,w_k\) be a fundamental system for \(L[y]=0\),
    with Wronskian matrix \(W\). Choose \function[functions] \(c_j\) satisfying
    \[
        W(x)\begin{pmatrix}c_1'\\ \vdots\\c_k'\end{pmatrix}
        =\begin{pmatrix}0\\ \vdots\\0\\ b(x)\end{pmatrix}.
    \]
    Then \(y_*=\sum_{j=1}^k c_jw_j\) is a particular solution of \(L[y]=b\).
    Such \(c_j\) exist on all of \(I\), by integrating the continuous vector
    \(W^{-1}(0,\ldots,0,b)^{\mathsf T}\).
\end{snippettheorem}

\begin{snippetproof}{diffeq-variation-constants-proof}{diffeq-variation-constants-theorem}{Variation of constants}
    The first \(k-1\) rows impose
    \(\sum_j c_j'w_j^{(r)}=0\) for \(0\leq r\leq k-2\).
    Successive differentiation therefore gives
    \[
        y_*^{(r)}=\sum_j c_jw_j^{(r)}\quad(0\leq r<k),\qquad
        y_*^{(k)}=\sum_j c_jw_j^{(k)}+\sum_j c_j'w_j^{(k-1)}.
    \]
    The last sum is \(b\), by the last row. Thus
    \(L[y_*]=\sum_j c_jL[w_j]+b=b\).
    In particular, for \(k=2\), the imposed equations are
    \[
        c_1'w_1+c_2'w_2=0,\qquad
        c_1'w_1'+c_2'w_2'=b.
    \]
    They make \(y_*'=c_1w_1'+c_2w_2'\) and produce exactly the desired
    forcing in the second derivative.
\end{snippetproof}

\begin{snippetexample}{diffeq-variation-first-order-example}{The first-order case}
    For \(y'+p(x)y=q(x)\), with \(p,q\in C(I)\), choose a primitive \(P'=p\).
    The homogeneous solution \(w=e^{-P}\) is never zero.
    Writing \(y_*=cw\) yields \(c'w=q\), so
    \[
        y=e^{-P(x)}\left(C+\int q(x)e^{P(x)}\,dx\right).
    \]
    This separates the homogeneous term \(Ce^{-P}\) from the particular
    term \(e^{-P}\int qe^P\) and agrees with the integrating-factor method.
\end{snippetexample}

\begin{snippetexample}{diffeq-variation-secant-example}{A nonpolynomial forcing term}
    Consider \(y''+y=1/\cos x\) on \(I=(-\pi/2,\pi/2)\).
    The characteristic equation \(\lambda^2+1=0\) gives the fundamental
    system \(w_1=\cos x,w_2=\sin x\).
    Set \(y_*=c_1(x)\cos x+c_2(x)\sin x\) and impose
    \[
        \begin{cases}
        c_1'\cos x+c_2'\sin x=0,\\
        -c_1'\sin x+c_2'\cos x=1/\cos x.
        \end{cases}
    \]
    The first equation gives \(c_1'=-c_2'\sin x/\cos x\).
    Substitution in the second yields
    \[
        c_2'\frac{\sin^2x+\cos^2x}{\cos x}=\frac1{\cos x},
        \qquad c_2'=1,\quad c_1'=-\tan x.
    \]
    As \(\cos x>0\) on \(I\), take \(c_2=x\), \(c_1=\ln(\cos x)\).
    The general solution is therefore
    \[
        y=k_1\cos x+k_2\sin x+\cos x\ln(\cos x)+x\sin x.
    \]
\end{snippetexample}
\end{document}
